\subsection{Overlapping-window corner cancellation}
\label{subsec:peps_torus_overlap_corner_cancellation}

\begin{definition}[Corner extension of a region insert]
    \label{def:peps_corner_extension}
    \lean{TNLean.PEPS.nestedThreeBlockGeometry}
    \lean{TNLean.PEPS.bareExtendInsert}
    \lean{TNLean.PEPS.extendInsert}
    \leanok
    Let $R\subseteq S$ be regions in a finite graph, and let $C$ be a region
    insert on $R$.  Write $\Phi_{R,S}$ for the coefficient obtained by
    contracting the added block $S\setminus R$ between the boundary of $R$ and
    the boundary of $S$.  The bare corner extension is
    \[
        \overline{\Ext}_{R\subseteq S}(C)_{\nu,\sigma}
        =
        \sum_{\mu\in \partial R}
        C_{\mu,\sigma|_R}\,
        \Phi_{R,S}\bigl(\mu,\sigma|_{S\setminus R},\nu\bigr).
    \]
    The normalized corner extension of $C$ from $R$ to $S$ is the region insert
    \[
        \Ext_{R\subseteq S}(C)_{\nu,\sigma}
        =
        P_{V\setminus S}^{-1}
        \sum_{\mu\in \partial R}
        C_{\mu,\sigma|_R}\,
        \Phi_{R,S}\bigl(\mu,\sigma|_{S\setminus R},\nu\bigr),
    \]
    where $P_{V\setminus S}$ is the product of bond dimensions over the edges
    not crossing the boundary of $V\setminus S$.  This is the open-boundary
    form of the extension used in the overlapping-window proof sketch of
    \cite[Section~4, proof of Theorem~3]{Molnar2018NormalPEPS}.
\end{definition}

\begin{theorem}[State preservation under corner extension]
    \label{thm:peps_corner_extension_state_preserves}
    \lean{TNLean.PEPS.deformedRegionState_extend}
    \leanok
    \uses{def:peps_corner_extension}
    If all bond dimensions are positive, then corner extension from
    $R\subseteq S$ preserves the deformed region state:
    \[
        \Psi_{A,R,C}
        =
        \Psi_{A,S,\Ext_{R\subseteq S}(C)}.
    \]
\end{theorem}
\begin{proof}
    \leanok
    Expanding the deformed state on $S$ separates the added block
    $S\setminus R$ from the complement $V\setminus S$.  The corner-extension
    sum contracts the added block first, giving the pointwise identity
    \[
        \Psi_{A,S,\Ext_{R\subseteq S}(C)}(\tau)
        =
        P_{V\setminus S}^{-1}\,P_{V\setminus S}\,
        \Psi_{A,R,C}(\tau)
        =
        \Psi_{A,R,C}(\tau).
    \]
\end{proof}

\begin{theorem}[Consecutive horizontal windows have equal corner extensions]
    \label{thm:peps_horizontal_consecutive_window_extend_eq}
    \lean{TNLean.PEPS.horizontalConsecutiveWindow_extend_eq}
    \leanok
    \uses{def:peps_corner_extension,
        thm:peps_corner_extension_state_preserves}
    Assume the one-orientation $L\times K$ window-injectivity hypotheses,
    union closure for injective regions, positive bond dimensions, and the
    bounds
    \[
        2\le L,\qquad 2\le K,\qquad 2L+1\le W,\qquad 2K+1\le H.
    \]
    Suppose two horizontally consecutive $L\times K$ windows carry inserts
    $C_1$ and $C_2$ whose deformed states agree.  Let $U$ be their
    $(L+1)\times K$ union.  Then their extensions to $U$ agree:
    \[
        \Ext_{W_1\subseteq U}(C_1)
        =
        \Ext_{W_2\subseteq U}(C_2).
    \]
\end{theorem}
\begin{proof}
    \leanok
    Theorem~\ref{thm:peps_corner_extension_state_preserves} rewrites each
    single-window deformed state as the deformed state on $U$ with the
    corresponding extended insert:
    \[
        \Psi_{A,U,\Ext_{W_1\subseteq U}(C_1)}
        =
        \Psi_{A,W_1,C_1}
        =
        \Psi_{A,W_2,C_2}
        =
        \Psi_{A,U,\Ext_{W_2\subseteq U}(C_2)}.
    \]
    Since $U$ is injective, equality of the two $U$-states implies equality
    of the two inserts:
    \[
        \Psi_{A,U,\Ext_{W_1\subseteq U}(C_1)}
        =
        \Psi_{A,U,\Ext_{W_2\subseteq U}(C_2)}
        \quad\Longrightarrow\quad
        \Ext_{W_1\subseteq U}(C_1)
        =
        \Ext_{W_2\subseteq U}(C_2).
    \]
\end{proof}

\paragraph{Uniform interior-bond products for staircase windows.}

For a PEPS tensor $A$ and a region $R$, write
\[
    P_A(R)=\prod_{e\notin\partial R} D_A(e)
\]
for the product of bond dimensions over the non-boundary edges of $R$.

\begin{theorem}[Interior-bond product under a graph isomorphism]%
    \label{thm:peps_regionInteriorBondProd_map}
    \lean{TNLean.PEPS.regionInteriorBondProd_map}
    \leanok
    Let $\phi:G\simeq G'$ be an isomorphism of finite graphs, and let $R$ be
    a region of $G$.  If $A'$ is a PEPS tensor on $G'$, then
    \[
        P_{A'}(\phi(R))
        =
        \prod_{e\notin\partial R} D_{A'}(\phi(e)).
    \]
\end{theorem}
\begin{proof}
    \leanok
    An edge of $\phi(R)$ is non-boundary exactly when its preimage is
    non-boundary for $R$:
    \[
        e'\notin\partial\phi(R)
        \Longleftrightarrow
        \phi^{-1}(e')\notin\partial R.
    \]
    Reindexing the product over non-boundary edges by the edge bijection
    $e\mapsto\phi(e)$ gives the displayed equality.
\end{proof}

\begin{theorem}[Translation invariance of the interior-bond product]%
    \label{thm:peps_regionInteriorBondProd_translate}
    \lean{TNLean.PEPS.regionInteriorBondProd_translate_of_translationInvariant}
    \leanok
    \uses{def:peps_torus_translation_invariant,
        thm:peps_regionInteriorBondProd_map}
    Let $A$ be translation invariant on a torus whose periods $W,H$ satisfy
    $1<W$ and $1<H$.  For every translation $\tau_{a,b}$ and every region $R$,
    \[
        P_A(\tau_{a,b}(R))=P_A(R).
    \]
\end{theorem}
\begin{proof}
    \leanok
    Theorem~\ref{thm:peps_regionInteriorBondProd_map} gives
    \[
        P_A(\tau_{a,b}(R))
        =
        \prod_{e\notin\partial R}D_A(\tau_{a,b}(e)).
    \]
    Translation invariance gives $D_A(\tau_{a,b}(e))=D_A(e)$ for every
    edge $e$, so the product is $P_A(R)$.
\end{proof}

\begin{theorem}[Interior-bond product of a cyclic rectangle is start-independent]%
    \label{thm:peps_regionInteriorBondProd_torusArcRectangle_eq}
    \lean{TNLean.PEPS.regionInteriorBondProd_torusArcRectangle_eq}
    \leanok
    \uses{thm:peps_regionInteriorBondProd_translate}
    Let $A$ be translation invariant on a torus whose periods $W,H$ satisfy
    $1<W$ and $1<H$.  If $Q(s;x,y)$ denotes the cyclic $x\times y$ rectangle
    starting at $s$, then for any starts $s,s'$,
    \[
        P_A(Q(s;x,y))=P_A(Q(s';x,y)).
    \]
\end{theorem}
\begin{proof}
    \leanok
    If $s=(s_x,s_y)$ and $s'=(s'_x,s'_y)$, then the translation by
    $(s_x-s'_x,s_y-s'_y)$ carries one rectangle to the other:
    \[
        \tau_{s_x-s'_x,s_y-s'_y}(Q(s';x,y))=Q(s;x,y).
    \]
    Apply Theorem~\ref{thm:peps_regionInteriorBondProd_translate}.
\end{proof}

\begin{theorem}[Uniform interior-bond product of the staircase windows]%
    \label{thm:peps_regionInteriorBondProd_staircaseWindow_eq}
    \lean{TNLean.PEPS.regionInteriorBondProd_staircaseWindow_eq}
    \leanok
    \uses{thm:peps_regionInteriorBondProd_torusArcRectangle_eq}
    Let $A$ be translation invariant on a torus whose periods $W,H$ satisfy
    $1<W$ and $1<H$.  If $W_j$ and $W_{j'}$ are two windows of the staircase
    based at the same corner $s$ with window dimensions $L\times K$, then
    \[
        P_A(W_j)=P_A(W_{j'}).
    \]
\end{theorem}
\begin{proof}
    \leanok
    Each staircase window is an $L\times K$ cyclic rectangle:
    \[
        W_j=Q(s_j;L,K),
        \qquad
        W_{j'}=Q(s_{j'};L,K).
    \]
    Theorem~\ref{thm:peps_regionInteriorBondProd_torusArcRectangle_eq}
    gives $P_A(Q(s_j;L,K))=P_A(Q(s_{j'};L,K))$.
\end{proof}

\begin{theorem}[Uniform horizontal per-step bond transport]%
    \label{thm:peps_horizontal_staircase_consecutive_window_bond_transport_uniform}
    \lean{TNLean.PEPS.horizontalStaircaseConsecutiveWindow_bondTransport_extend_eq_uniform}
    \leanok
    \uses{def:peps_corner_extension,
        thm:peps_regionInteriorBondProd_staircaseWindow_eq}
    Let $B$ be translation invariant on a torus whose periods $W,H$ satisfy
    $1<W$ and $1<H$.  Assume the one-orientation $L\times K$
    window-injectivity hypotheses, union closure for injective regions,
    positive bond dimensions, and the bounds
    \[
        2\le L,\qquad 2\le K,\qquad 2L+1\le W,\qquad 2K+1\le H.
    \]
    Let $j<L$, let $U_j=W_j\cup W_{j+1}$ be the horizontal consecutive-window
    union, and let $g$ be a boundary edge of both $W_j$ and $W_{j+1}$.  If the
    matrices obtained from $M_1$ and $M_2$ by the orientation convention at
    $g$ agree, then, with $P_j=P_B(W_j)$,
    \[
        \operatorname{Ext}_{W_j\subseteq U_j}
            \bigl(P_j\,I_{B,W_j,g,M_1}\bigr)
        =
        \operatorname{Ext}_{W_{j+1}\subseteq U_j}
            \bigl(P_j\,I_{B,W_{j+1},g,M_2}\bigr).
    \]
\end{theorem}
\begin{proof}
    \leanok
    The consecutive-window transport gives the cross-rescaled equality
    \[
        \operatorname{Ext}_{W_j\subseteq U_j}
            \bigl(P_B(W_{j+1})\,I_{B,W_j,g,M_1}\bigr)
        =
        \operatorname{Ext}_{W_{j+1}\subseteq U_j}
            \bigl(P_B(W_j)\,I_{B,W_{j+1},g,M_2}\bigr).
    \]
    By Theorem~\ref{thm:peps_regionInteriorBondProd_staircaseWindow_eq},
    \[
        P_B(W_{j+1})=P_B(W_j)=P_j.
    \]
    Substituting this equality gives the displayed common-scalar form.
\end{proof}

\begin{theorem}[Uniform vertical per-step bond transport]%
    \label{thm:peps_vertical_staircase_consecutive_window_bond_transport_uniform}
    \lean{TNLean.PEPS.verticalStaircaseConsecutiveWindow_bondTransport_extend_eq_uniform}
    \leanok
    \uses{def:peps_corner_extension,
        thm:peps_regionInteriorBondProd_staircaseWindow_eq}
    Let $B$ be translation invariant on a torus whose periods $W,H$ satisfy
    $1<W$ and $1<H$.  Assume the one-orientation $L\times K$
    window-injectivity hypotheses, union closure for injective regions,
    positive bond dimensions, and the bounds
    \[
        2\le L,\qquad 2\le K,\qquad 2L+1\le W,\qquad 2K+1\le H.
    \]
    Let $L\le j$ and $j+1<L+K$, let $U_j=W_j\cup W_{j+1}$ be the vertical
    consecutive-window union, and let $g$ be a boundary edge of both $W_j$ and
    $W_{j+1}$.  If the matrices obtained from $M_1$ and $M_2$ by the
    orientation convention at $g$ agree, then, with $P_j=P_B(W_j)$,
    \[
        \operatorname{Ext}_{W_j\subseteq U_j}
            \bigl(P_j\,I_{B,W_j,g,M_1}\bigr)
        =
        \operatorname{Ext}_{W_{j+1}\subseteq U_j}
            \bigl(P_j\,I_{B,W_{j+1},g,M_2}\bigr).
    \]
\end{theorem}
\begin{proof}
    \leanok
    The descending-arm transport gives the same cross-rescaled equality as in
    the horizontal case, with $U_j$ now the $L\times(K+1)$ union:
    \[
        \operatorname{Ext}_{W_j\subseteq U_j}
            \bigl(P_B(W_{j+1})\,I_{B,W_j,g,M_1}\bigr)
        =
        \operatorname{Ext}_{W_{j+1}\subseteq U_j}
            \bigl(P_B(W_j)\,I_{B,W_{j+1},g,M_2}\bigr).
    \]
    The uniformity theorem gives $P_B(W_{j+1})=P_B(W_j)=P_j$.  Substituting
    into the displayed cross-rescaled form gives
    \[
        \operatorname{Ext}_{W_j\subseteq U_j}
            \bigl(P_j\,I_{B,W_j,g,M_1}\bigr)
        =
        \operatorname{Ext}_{W_{j+1}\subseteq U_j}
            \bigl(P_j\,I_{B,W_{j+1},g,M_2}\bigr).
    \]
\end{proof}


\begin{theorem}[Patch-extended equality of the staircase end windows]
    \label{thm:peps_horizontal_staircase_patch_extend_eq}
    \lean{TNLean.PEPS.horizontalStaircase_patch_extend_eq}
    \leanok
    \uses{def:peps_horizontal_staircase_patch,
        thm:peps_corner_extension_state_preserves}
    Assume positive bond dimensions and the bounds
    \[
        0<L,\qquad 0<K,\qquad 2L\le W,\qquad 2K\le H.
    \]
    If the deformed states of the two end-window inserts agree, then their
    extensions to the staircase patch $P$ have equal deformed states:
    \[
        \Psi_{A,P,\Ext_{W_{\mathrm{right}}\subseteq P}(C_0)}
        =
        \Psi_{A,P,\Ext_{W_{\mathrm{left}}\subseteq P}(C_{\mathrm{end}})}.
    \]
\end{theorem}
\begin{proof}
    \leanok
    Applying Theorem~\ref{thm:peps_corner_extension_state_preserves} to each
    end window gives
    \[
        \Psi_{A,P,\Ext_{W_{\mathrm{right}}\subseteq P}(C_0)}
        =
        \Psi_{A,W_{\mathrm{right}},C_0}
        =
        \Psi_{A,W_{\mathrm{left}},C_{\mathrm{end}}}
        =
        \Psi_{A,P,\Ext_{W_{\mathrm{left}}\subseteq P}(C_{\mathrm{end}})},
    \]
    where the middle equality is the hypothesis.
\end{proof}

% The end-pair stripping under full complement injectivity has been retired: at the
% corollary's minimal size $V\setminus S$ is not blocked-tensor injective, so its hypothesis
% is not derivable.  The faithful Step~3 is the open-boundary shared-corner cancellation
% \ref{thm:peps_staircase_pair_insert_eq_open}, which inverts only the injective completed
% corner $Q$ and never $V\setminus S$.

\begin{lemma}[Two-block merge collapse]
    \label{lem:peps_two_block_merge_collapse}
    \lean{TNLean.PEPS.mergeCollapse2}
    \leanok
    Let $B_1,K_1\subseteq V\setminus T$ and $B_2,H_2\subseteq T$.  Fix
    boundary data $c_1$ on $V\setminus K_1$ and $c_2$ on $H_2$, and physical
    labels $\sigma_1$ on $B_1$ and $\sigma_2$ on $B_2$.  Set
    \[
        F_1(\xi)=\prod_{x\in B_1}A_x(\xi|_{\partial x},\sigma_1(x)),
        \qquad
        F_2(\xi)=\prod_{x\in B_2}A_x(\xi|_{\partial x},\sigma_2(x)).
    \]
    Then the
    boundary-compatible double sum collapses to the multiplicity of the
    interior bonds of $T$ times the single merged sum:
    \[
        \sum_{\substack{
            \zeta_2|_{\partial(V\setminus K_1)}=c_1,\;
            \zeta_1|_{\partial H_2}=c_2\\
            \zeta_1|_{\partial T}=\zeta_2|_{\partial T}}}
        F_1(\zeta_2)F_2(\zeta_1)
        =
        P_T
        \sum_{\substack{
            \eta|_{\partial(V\setminus K_1)}=c_1,\;
            \eta|_{\partial H_2}=c_2}}
        F_1(\eta)F_2(\eta).
    \]
\end{lemma}
\begin{proof}
    \leanok
    Each merged configuration has exactly $P_T$ preimages, indexed by the
    virtual bonds strictly inside $T$.  Since $B_1$ and $K_1$ lie outside
    $T$, while $B_2$ and $H_2$ lie inside $T$, the two products and the two
    boundary conditions are unchanged under the merge.
\end{proof}

\begin{theorem}[Composition of corner-extension coefficient kernels]
    \label{thm:peps_corner_extension_coefficient_composition}
    \lean{TNLean.PEPS.threeBlockBlueCoeff_comp}
    \leanok
    \uses{def:peps_corner_extension, lem:peps_two_block_merge_collapse}
    For nested regions $R\subseteq S\subseteq P$, the coefficient kernels
    satisfy
    \[
        \sum_{\nu\in\partial S}
        \Phi_{R,S}\bigl(\mu,\sigma|_{S\setminus R},\nu\bigr)
        \Phi_{S,P}\bigl(\nu,\sigma|_{P\setminus S},\omega\bigr)
        =
        P_{V\setminus S}\,
        \Phi_{R,P}\bigl(\mu,\sigma|_{P\setminus R},\omega\bigr).
    \]
    This is the two-block merge collapse in the proof of
    \cite[Section~4, proof of Theorem~3]{Molnar2018NormalPEPS}.
\end{theorem}

\begin{theorem}[Linearity and transitivity of corner extension]
    \label{thm:peps_corner_extension_linear_transitive}
    \lean{TNLean.PEPS.bareExtendInsert_const_smul}
    \lean{TNLean.PEPS.extendInsert_const_smul}
    \lean{TNLean.PEPS.extendInsert_add}
    \lean{TNLean.PEPS.extendInsert_sub}
    \lean{TNLean.PEPS.bareExtendInsert_trans}
    \lean{TNLean.PEPS.extendInsert_trans}
    \leanok
    \uses{def:peps_corner_extension,
        thm:peps_corner_extension_coefficient_composition}
    The corner extension is linear in the inserted boundary coefficients.  For
    nested regions $R\subseteq S\subseteq P$, the bare extension satisfies
    \[
        \overline{\Ext}_{S\subseteq P}
        \bigl(\overline{\Ext}_{R\subseteq S}(C)\bigr)
        =
        P_{V\setminus S}\,
        \overline{\Ext}_{R\subseteq P}(C)
    \]
    without a positivity assumption.  If all bond dimensions are positive, the
    normalized corner extension is transitive:
    \[
        \Ext_{S\subseteq P}
        \bigl(\Ext_{R\subseteq S}(C)\bigr)
        =
        \Ext_{R\subseteq P}(C).
    \]
\end{theorem}
\begin{proof}
    \leanok
    The scalar and additive laws follow from the defining sum.  For
    transitivity, substitute
    Theorem~\ref{thm:peps_corner_extension_coefficient_composition} into the
    defining sum.  This gives the bare formula, and the factors
    $P_{V\setminus S}$ cancel in the normalized extension.
\end{proof}

\begin{definition}[Assembly of subregion physical configurations]
    \label{def:peps_assemble_subregion_physical_config}
    \lean{TNLean.PEPS.assembleSubRegionσ}
    \leanok
    If $R\subseteq S$, a physical configuration on $R$ and a physical
    configuration on $S\setminus R$ assemble to a physical configuration on
    $S$:
    \[
        \operatorname{As}_{R,S}(\sigma_R,\sigma_Q)(v)
        =
        \begin{cases}
        \sigma_R(v), & v\in R,\\
        \sigma_Q(v), & v\in S\setminus R.
        \end{cases}
    \]
\end{definition}

\begin{theorem}[Restrictions of assembled subregion configurations]
    \label{thm:peps_assemble_subregion_physical_config_restricts}
    \lean{TNLean.PEPS.restrictSubRegionσ_assembleSubRegionσ}
    \lean{TNLean.PEPS.restrictSubRegionσ_sdiff_assembleSubRegionσ}
    \leanok
    \uses{def:peps_assemble_subregion_physical_config}
    The two restrictions recover the original configurations:
    \[
        \operatorname{As}_{R,S}(\sigma_R,\sigma_Q)|_R=\sigma_R,
        \qquad
        \operatorname{As}_{R,S}(\sigma_R,\sigma_Q)|_{S\setminus R}=\sigma_Q.
    \]
\end{theorem}

\begin{theorem}[Kernel triviality of corner extension from the added block]
    \label{thm:peps_extend_insert_kernel_trivial_added_injective}
    \lean{TNLean.PEPS.extendInsert_kernel_trivial_of_addedInjective}
    \leanok
    \uses{def:peps_corner_extension,
        def:peps_assemble_subregion_physical_config,
        thm:peps_assemble_subregion_physical_config_restricts}
    Let $R\subseteq S$ be finite regions.  Suppose that all bond dimensions are
    positive and that the added block $Q=S\setminus R$ is blocked-tensor
    injective.  If a region insert $C$ on $R$ has zero corner extension,
    \[
        \Ext_{R\subseteq S}(C)=0,
    \]
    then $C=0$.
\end{theorem}
\begin{proof}
    \leanok
    Since all bond dimensions are positive, $P_{V\setminus S}\ne 0$.
    Thus $\Ext_{R\subseteq S}(C)=0$ is equivalent to
    $\overline{\Ext}_{R\subseteq S}(C)=0$.  Fix the physical
    configuration $\sigma_R$ on $R$.  Let
    \[
        \rho:\partial R \simeq \partial(V\setminus R)
        \qquad\text{and}\qquad
        \tau:\partial S \simeq \partial(V\setminus S)
    \]
    be the complement-boundary bijections, and define
    \[
        c_\eta=C_{\rho^{-1}(\eta),\sigma_R}
        \qquad(\eta\in \partial(V\setminus R)).
    \]
    Evaluating the bare zero extension at
    $\operatorname{As}_{R,S}(\sigma_R,\sigma_Q)$ gives, for every added-block
    configuration $\sigma_Q$ and every boundary value
    $\omega\in\partial(V\setminus S)$,
    \[
        \sum_{\eta\in\partial(V\setminus R)}
        c_\eta\,\Phi_{R,S}(\rho^{-1}(\eta),\sigma_Q,\tau^{-1}(\omega))=0.
    \]
    For each outer boundary value $\omega$ and each boundary value $\beta$ of $Q$,
    define
    \[
        f_\omega(\beta)
        =
        \sum_{\eta\in\partial(V\setminus R)}
        c_\eta\,
        \mathbf 1_{\{\exists q:\,
            \partial_{V\setminus R}q=\eta,\,
            \partial_{V\setminus S}q=\omega,\,
            \partial_Qq=\beta\}}.
    \]
    Let $\Gamma$ be the product of bond dimensions on the crossing edges
    separated when the $Q$-block is read first.  Unfolding $\Phi_{R,S}$ as a sum
    over the blocked weights of $Q$ gives, for every $\sigma_Q$,
    \[
    \begin{aligned}
        (T_Q f_\omega)(\sigma_Q)
        &=
        \sum_{\beta\in\partial Q} f_\omega(\beta)\,W_Q(\beta,\sigma_Q) \\
        &=
        \sum_{\beta\in\partial Q}
        \left(
        \sum_{\eta\in\partial(V\setminus R)}
        c_\eta\,
        \mathbf 1_{\{\exists q:\,
            \partial_{V\setminus R}q=\eta,\,
            \partial_{V\setminus S}q=\omega,\,
            \partial_Qq=\beta\}}
        \right) W_Q(\beta,\sigma_Q) \\
        &=
        \Gamma
        \sum_{\eta\in\partial(V\setminus R)}
        c_\eta\,\Phi_{R,S}(\rho^{-1}(\eta),\sigma_Q,\tau^{-1}(\omega))
        =
        0.
    \end{aligned}
    \]
    Hence $T_Q(f_\omega)=0$.  Injectivity of the blocked tensor of $Q$ gives
    \[
        T_Q(f_\omega)=0 \Longrightarrow f_\omega=0.
    \]
    Now fix $\mu\in\partial R$.  Positivity of bond dimensions gives a global
    virtual configuration $q$ whose residual boundary on $V\setminus R$ is
    $\rho(\mu)$.  Set
    \[
        \omega=\partial_{V\setminus S}q,\qquad
        \beta=\partial_Q q.
    \]
    Evaluating $f_\omega=0$ at $\beta$, the residual-boundary uniqueness identity
    leaves only the term $\eta=\rho(\mu)$:
    \[
        0
        =
        f_\omega(\beta)
        =
        c_{\rho(\mu)}\cdot 1
        =
        C_{\mu,\sigma_R}.
    \]
    Since $\sigma_R$ and $\mu$ were arbitrary, $C=0$.
\end{proof}

\begin{theorem}[Cancellation of a shared injective corner]
    \label{thm:peps_extend_insert_cancel_added_injective}
    \lean{TNLean.PEPS.extendInsert_cancel_addedInjective}
    \leanok
    \uses{thm:peps_corner_extension_linear_transitive,
        thm:peps_extend_insert_kernel_trivial_added_injective}
    Let $R\subseteq S$, and suppose that all bond dimensions are positive and
    that $S\setminus R$ is blocked-tensor injective.  If two inserts on $R$
    have the same corner extension to $S$,
    \[
        \Ext_{R\subseteq S}(C_1)
        =
        \Ext_{R\subseteq S}(C_2),
    \]
    then $C_1=C_2$.
\end{theorem}
\begin{proof}
    \leanok
    By linearity,
    \[
        \Ext_{R\subseteq S}(C_1-C_2)=0.
    \]
    Theorem~\ref{thm:peps_extend_insert_kernel_trivial_added_injective}
    applied to the added block $S\setminus R$ gives $C_1-C_2=0$, hence
    $C_1=C_2$.
\end{proof}

\begin{theorem}[Staircase-pair cancellation on the torus]
    \label{thm:peps_staircase_pair_cancel}
    \lean{TNLean.PEPS.staircasePair_cancel}
    \leanok
    \uses{thm:peps_extend_insert_cancel_added_injective}
    Let the torus have width $W$ and height $H$, with $1<W$ and $1<H$.
    Let $S$ be the staircase end pair around a horizontal edge, with seam
    dimensions $L$ and $K$, and let $Q$ be the completed corner rectangle.
    Assume the size bounds
    \[
        2L\le W,\qquad 2K\le H,
    \]
    that $Q$ is blocked-tensor injective, and that all bond dimensions are
    positive.
    If two inserts $X$ and $Y$ on $S$ have equal extensions to $S\cup Q$,
    \[
        \Ext_{S\subseteq S\cup Q}(X)
        =
        \Ext_{S\subseteq S\cup Q}(Y),
    \]
    then $X=Y$.  Thus the injective completed corner can be cancelled from an
    open-boundary equality, without assuming injectivity of the torus
    complement of $S$.
\end{theorem}
\begin{proof}
    \leanok
    The size bounds imply that the staircase end pair is disjoint from the
    completed corner, and therefore
    \[
        (S\cup Q)\setminus S=Q.
    \]
    The claimed equality is exactly
    Theorem~\ref{thm:peps_extend_insert_cancel_added_injective} applied to the
    inclusion $S\subseteq S\cup Q$.
\end{proof}
